% P1-2: intro (systems problem + 3 properties) + compact method + F1.

\section{Introduction}
\label{sec:intro}

For an agent to hold persistent memory of its environment, its internal
state cannot grow unboundedly with time. Standard retrieval-augmented
systems embed every observation into a dense vector bank that grows
linearly with recording duration, untenable for continuous multi-session
memory on-device. Improving the long-context capacity of the models
themselves~\cite{arora2024linear,chen2024vljepa} does not address this
systems bottleneck: how to structure a dynamic spatial representation
supporting test-time memory updates without unbounded state explosion.

Three properties define a viable substrate for multi-session persistent
spatial memory. \emph{Memory-budgeted}: state bounded independently of
sequence length, scaling by physical area explored. \emph{Dynamic}:
test-time evolution, letting old information decay while maintaining
unique-event cardinality across revisits. \emph{Streaming}: in-place,
online updates. We present \textbf{Probabilistic Spatial Memory}
(\psm{}), which meets all three, and demonstrate the multi-session regime
it targets on 102.9\,h of metropolitan driving with genuine same-place
revisits---the setting a single-pass egocentric corpus cannot exercise.

\section{Probabilistic Spatial Memory}
\label{sec:method}

\psm{} is a streaming spatial-temporal index that ingests
$(t, \mathrm{lat}, \mathrm{lng}, \mathbf{e})$ tuples---timestamp,
geographic position, and a $d$-dimensional embedding---and supports
bounded-memory top-$K$ similarity search at any time
(Fig.~\ref{fig:f1-schematic}). Three layered probabilistic data structures
over an H3~\cite{brodsky2018h3} hexagonal tessellation enforce the memory
bound.

\paragraph{H3 tessellation.}
Each tuple maps to its containing H3 cell at resolution~$r$ (default~12,
${\sim}9$\,m edges; the street-scale results below use $r{=}10$,
${\sim}66$\,m). Cell IDs are global---comparable across recordings,
sessions, and devices---which is what makes cross-session and
cross-device memory well-defined.

\paragraph{Time-decayed HLL ring buffer.}
Each cell keeps a ring buffer of $C$ HyperLogLog
sketches~\cite{flajolet2007hyperloglog} rotating on a fixed time-window
cadence (default 30\,s). Ingest hashes the embedding into the current
sketch; when the window elapses the pointer advances and the oldest sketch
expires, enforcing bounded-time retention. The sketches estimate
unique-event cardinality per cell and expose the cell's temporal envelope
$(t_{\min}, t_{\max})$. At precision $p{=}10$ each sketch is $1\,$KiB, so
a cell's ring is $C\cdot1\,\mathrm{KiB} = 60\,$KiB, independent of
recording length.

\paragraph{Reservoir-sampled exemplars.}
For cosine scoring we store embeddings via Algorithm~R reservoir
sampling~\cite{vitter1985reservoir} per cell (capacity $R{=}128$): after
$n > R$ frames a cell holds a uniform $R$-sized sample. Brute-force grows
$\Theta(N)$; \psm{} grows $\Theta(R\cdot|\mathrm{cells}|)$, bounded by the
physical area covered, not recording length---the design's central
property.

\paragraph{Top-$K$ query.}
A query embeds the request and, for each cell in an optional
$k$-ring neighborhood, returns its top-$\min(\caponek, R)$ exemplars by
cosine similarity, then merges and truncates to $K$. \caponek{}$=1$
enforces strict spatial diversity (one slot per cell); \caponek{}$=K$
recovers unconstrained top-$K$. Ingest is $O(1)$ amortized (H3 lookup, HLL
hash, reservoir update); queries run at any time after the first ingest.
For look-back QA the ranked candidates feed a frontier-MLLM reranker
(full pipeline in the extended paper).

% F1-compact — PSM schematic for the 4pp abstract (single row, ~1/4 col).
% Shows the ingest->per-cell-state->query core only; the MLLM-rerank
% pipeline is deferred to the full paper. Reuses the visual vocabulary of
% figures/f1_architecture.tex (hex cells, HLL ring, reservoir slots).
% Requires \usetikzlibrary{positioning,arrows.meta,fit,calc,
% shapes.geometric,backgrounds} (loaded in main_4pp.tex).
\begin{figure}[t]
\centering
\begin{tikzpicture}[
    >={Stealth[length=1.4mm,width=1.1mm]},
    font=\scriptsize,
    box/.style={draw, rounded corners=1pt, align=center,
                inner sep=1.5pt, minimum height=4.5mm, fill=white},
    iobox/.style={box, fill=gray!8},
    hex/.style={draw, regular polygon, regular polygon sides=6,
                minimum size=4.4mm, inner sep=0pt, rotate=30, fill=gray!5},
    hexhi/.style={hex, fill=blue!12, draw=blue!60!black, thick},
    hllbox/.style={draw, minimum width=1.6mm, minimum height=1.6mm,
                   inner sep=0pt, fill=orange!25, draw=orange!70!black},
    hllold/.style={hllbox, fill=orange!8, draw=orange!30},
    resvslot/.style={draw, minimum width=1.6mm, minimum height=1.6mm,
                     inner sep=0pt, fill=blue!15, draw=blue!60!black},
    flow/.style={->, semithick, gray!60!black},
    smallflow/.style={->, thin, gray!60!black},
]

% ---- streaming input ----
\node[iobox] (in) {stream\\$(\mathbf{e},{\rm lat,lng},t)$};

% ---- H3 hex cluster ----
\coordinate (hexcenter) at ($(in.east) + (11mm, 0mm)$);
\node[hex]   (h0) at ($(hexcenter) + (-3.8mm, 2.2mm)$) {};
\node[hex]   (h1) at ($(hexcenter) + ( 0mm,   4.4mm)$) {};
\node[hex]   (h2) at ($(hexcenter) + ( 3.8mm, 2.2mm)$) {};
\node[hexhi] (hC) at (hexcenter) {};
\node[hex]   (h4) at ($(hexcenter) + (-3.8mm,-2.2mm)$) {};
\node[hex]   (h5) at ($(hexcenter) + ( 0mm,  -4.4mm)$) {};
\node[hex]   (h6) at ($(hexcenter) + ( 3.8mm,-2.2mm)$) {};
\node[font=\tiny, below=4.6mm of hC, gray!60!black] {H3 cells};

% ---- per-cell state pop-out ----
\coordinate (cellAnchor) at ($(hexcenter) + (12mm, 1mm)$);
% HLL ring (top row) with label ABOVE it
\node[hllold] (hll1) at ($(cellAnchor) + (-3mm, 3.0mm)$) {};
\node[hllold] (hll2) at ($(cellAnchor) + (-1mm, 3.0mm)$) {};
\node[hllbox] (hll3) at ($(cellAnchor) + ( 1mm, 3.0mm)$) {};
\node[hllbox] (hll4) at ($(cellAnchor) + ( 3mm, 3.0mm)$) {};
\node[font=\tiny, above=0.3mm of hll2, anchor=south, xshift=1mm]
     {HLL ring \textit{(decay)}};
% Reservoir (bottom row) with label BELOW it
\node[resvslot] (r1) at ($(cellAnchor) + (-3mm, 0.6mm)$) {};
\node[resvslot] (r2) at ($(cellAnchor) + (-1mm, 0.6mm)$) {};
\node[resvslot] (r3) at ($(cellAnchor) + ( 1mm, 0.6mm)$) {};
\node[resvslot] (r4) at ($(cellAnchor) + ( 3mm, 0.6mm)$) {};
\node[font=\tiny, below=0.3mm of r2, anchor=north, xshift=1mm] (rlab)
     {$R$ exemplars \textit{(bounded)}};
\draw[smallflow, dashed, blue!50!black]
     (hC.east) -- ($(cellAnchor) + (-5mm, 1.8mm)$);

% ---- query ----
\node[box, right=24mm of cellAnchor, align=center, fill=orange!10,
      draw=orange!70!black] (topk)
     {top-$K$ cosine\\{\tiny\caponek}};
\node[iobox, right=2.5mm of topk] (out)
     {$(t,{\rm cell})$\\$\times K$};

% ---- flows ----
\draw[flow] (in.east) -- (hC.west);
% reservoir -> top-K: exit right from the reservoir row, clear of labels
\draw[flow, dashed, blue!50!black]
     ($(r4.east)+(0.8mm,0)$) -- (topk.west);
\draw[flow] (topk.east) -- (out.west);

\end{tikzpicture}
\caption{\textbf{\psm{} schematic.} Streaming
$(\mathbf{e},\mathrm{lat},\mathrm{lng},t)$ tuples are deposited into H3
cells; each cell keeps a bounded ring of HLL sketches (orange, fading with
age) and $R$ reservoir-sampled embeddings (blue). A query cosine-searches
the reservoirs under the \caponek{} gate and returns $K$ cell-tagged
candidates. State is bounded per cell, independent of recording length.}
\label{fig:f1-schematic}
\end{figure}

